\tocsection{\textbf{9/9/26 Lecture 6: }Properties of the Lebesgue Measure}

\begin{lemma}[\textbf{Proposition 6.1}]
    There exists a set in $\cal{L}$ that is not in $\cal{B}$.
\end{lemma}
\begin{proof}{Proof:}
    First, we assume the fact that $\abs{\cal{B}} = \abs{\R}$.\\

    Now, let $C$ be the middle thirds Cantor set. We know that $\abs{C} = \abs{\R}$. Now, every subset of $C$ is Lebesgue measureable, since $C$ is a null set of $\lambda^{*}$, so $\cal{P}(C) \subseteq \cal{L}$. However, by Cantor's theorem, we see that $\abs{\cal{P}(C)} > \abs{C} = \abs{\R}$, so
    \[ \abs{\cal{L}} \geq \abs{\cal{P}(C)} > \abs{C} = \abs{\R} = \abs{\cal{B}} \]
    so as $\abs{\cal{B}} < \abs{\cal{L}}$, the desired result holds. 

    $\qed$
\end{proof}

It turns out that this construction entirely highlights the issue with $\cal{B}$: it doesn't contain all the null sets of $\lambda$, which we will soon see.

\begin{definition}[\textbf{Definition 6.2: }Complete Metric Spaces]
    A measure space $(X, \cal{A}, \mu)$ is \emph{complete} if $B \in \cal{A}$ and $\mu(B) = 0$ implies that for all $A \subseteq B$, $A \in \cal{A}$. 
\end{definition}

Notice that we can define a measure space $(\R, \cal{B}, \lambda)$ to be a measure space by restricting the Lebesgue measure to $\cal{B}$, but this is not a complete measure space.

\begin{lemma}[\textbf{Proposition 6.3}]
    Let $A \subseteq [0, 1]$ be $\lambda$-measurable. Then,
    \begin{enumerate}
        \item For any $\epsilon > 0$, there exists an open set $\cal{U}$ such that $A \subseteq \cal{U}$ and $\lambda(\cal{U} \backslash A) < \epsilon$.
        \item For any $\epsilon > 0$, there exists a closed set $\cal{C}$ such that $\cal{C} \subseteq A$ and $\lambda{A \backslash \cal{C}} < \epsilon$. 
        \item There are open sets
        \[ \cal{U}_1 \supseteq \cal{U}_2 \supseteq \cal{U}_3 \supseteq \cdots \supseteq A \]
        so that $\lambda\left(\bigcap_{i = 1}^{\infty} \cal{U}_i \backslash A\right) = 0$. 
        \item There are closed sets
        \[ \cal{C}_1 \subseteq \cal{C}_2 \subseteq \cdots \subseteq A \]
        such that $\lambda\left(A \backslash \bigcup_{i = 1}^{\infty} E_{n}\right) = 0$. 
    \end{enumerate}
\end{lemma}
\begin{proof}{Proof:}
    
    \textbf{1)} From the definition of $\lambda$, one can find an open set $\cal{U}$ such that $A \subseteq \cal{U}$ and $\lambda(U) < \lambda(A) + \epsilon$. Then, as $\lambda(\cal{U} \backslash A) = \lambda(\cal{U}) - \lambda(A)$ as $A \subseteq \cal{U}$, we see that $\lambda(\cal{U} \backslash A) < \epsilon$.\\

    \textbf{2)} Now, consider the set $[0, 1] \backslash A$, and using $1$ find an open set $\cal{O} \supseteq [0, 1] \backslash A$ such that $\lambda(\cal{O} \backslash ([0, 1] \backslash A)) < \epsilon$. Then, as 
    \begin{align*}
        \lambda(\cal{O} \backslash ([0, 1] \backslash A)) &= \lambda(\cal{O}) - \left(1 - \lambda(A)\right)\\
        &= \lambda(A) - (1 - \lambda(\cal{O})) \\
        &= \lambda(A) - \lambda([0, 1] \backslash \cal{O})
    \end{align*}
    and we see that $[0, 1] \backslash \cal{O}$ is a closed set contained in $A$, we see that this last expression is equal to $\lambda(A \backslash ([0, 1] \backslash \cal{O})) < \epsilon$.\\

    \textbf{3)} By $1$, there exists open sets $V_n \supseteq A$ such that $\lambda(V_n \backslash A) < 1/n$. Let $U_n = \bigcap_{k = 1}^{n} V_n$. Then, we see that
    \[ U_1 \supseteq U_2 \supseteq \cdots \supseteq A \]
    the $U_n$ are open, and $\lambda\left(\bigcap_{n = 1}^{\infty} V_n \backslash A\right) = \lim_{n \to \infty} \lambda(U_n \backslash A)$, which equals zero as $\lambda(U_n \backslash A) < 1/n$ for each $n$ (as $U_n \subseteq V_n$).\\

    \textbf{4)} Similar proof as \textbf{3)}. 

    $\qed$
\end{proof}

\begin{lemma}[\textbf{Corollary 6.4}]
    $\cal{L}$ is generated by $\cal{B}$ and the null sets of $\lambda$. 
\end{lemma}
\begin{proof}{Proof:}
    Let $A \in \cal{L}$. From the proposition above, there exists a $G_n = \cap_{i = 1}^{\infty} \cal{U}_i \in B$ so that $G_n \supseteq A_n = A \cap [n, n + 1]$ and $\lambda(G_n \backslash A_n) = 0$. Thus, we see that $\bigcup_{i = 1}^{\infty} G_n \in \cal{B}$ and
    \[ A = \left(\bigcup_{i = 1}^{\infty} G_n\right) \backslash \left(\bigcup_{i = 1}^P{\infty} G_n \backslash A_n\right) \]
    so $A$ can be written as the union, complement, and intersection of Borel and null sets. 

    $\qed$
\end{proof}

\begin{definition}[\textbf{Definition 6.5:} $G_{\delta}$ and $F_{\sigma}$ Sets]
    A countable intersection of open sets is a $G_{\delta}$ set. A countable union of closed sets is a $F_{\sigma}$ set. 
\end{definition}

\subsection*{The Lebesgue-Stieljes Measure}

Let $\cal{C}$ be a collection of intervals of the form $(a, b]$ and let $\alpha : \R \to \R$ be increasing and right-continuous. Then, define $\ell((a, b]) = \alpha(b) - \alpha(a)$. Now, for $E \subseteq \R$, define
\[ m^{*}(E) = \inf\{\sum_{n = 1}^{\infty} \ell(A_n) \col A_n \in \cal{C} \text{ and } E \subseteq \bigcup_{n = 1}^{\infty} A_n\} \]
For example, if $\alpha = x$, then we get the Lebesgue outer measure, and if $\alpha$ is zero for all $x < 0$ and $1$ for all $x \geq 0$, we see that $m^{*}$ is the point mass outer measure. 

\begin{lemma}[\textbf{Proposition 6.6}]
    The function $m^{*} : \cal{P} \to [0, \infty]$ is an outer measure. Let $m$ be the measure defined by the restriction of $m^{*}$ to $m^{*}$-measurable sets. Then, every Borel set is $m^{*}$-measurable and $m^{*}((a, b]) = \ell(a, b]$. 
\end{lemma}
\begin{proof}{Proof:}
    This is the same proof for $\lambda^{*}$, but we replace $\abs{(a, b)} = b - a$ with $\alpha(b) - \alpha(a)$. 
\end{proof}